\section{Kiến trúc hệ thống và luồng thực thi}
\subsection{Kiến trúc}
Nhóm em chia chương trình thành các phần tương đối độc lập. Graph contract và các thuật toán nằm ở phần core, không phụ thuộc trực tiếp vào FastAPI. ScenarioCostEngine chịu trách nhiệm tính chi phí edge và heuristic lower-bound. SearchService chạy thuật toán tìm kiếm và lưu trace. PairwiseMatrix tạo ma trận chi phí giữa các điểm, sau đó TourOptimizerService dùng ma trận này cho Held-Karp hoặc Nearest Neighbor. FastAPI cung cấp API, còn React/MapLibre hiển thị bản đồ và kết quả.

\begin{figure}[htbp]
\centering
\begin{tikzpicture}[
  box/.style={draw, rounded corners, minimum width=2.55cm, minimum height=0.78cm, align=center, fill=blue!7},
  core/.style={draw, rounded corners, minimum width=2.55cm, minimum height=0.78cm, align=center, fill=green!10},
  arrow/.style={->, thick, >=stealth}
]
\node[box] (ui) {React + MapLibre\\GUI};
\node[box, right=1.1cm of ui] (api) {FastAPI\\API};
\node[core, right=1.1cm of api] (search) {SearchService\\UCS / A* / trace};
\node[core, below=0.9cm of search] (scenario) {ScenarioCostEngine\\Graph contract};
\node[core, below=0.9cm of api] (tour) {PairwiseMatrix\\TourOptimizerService};
\draw[arrow] (ui) -- (api);
\draw[arrow] (api) -- (search);
\draw[arrow] (api) -- (tour);
\draw[arrow] (search) -- (scenario);
\draw[arrow] (tour) -- (search);
\end{tikzpicture}
\caption{Sơ đồ các thành phần chính của FloodRoute HCMC.}
\label{fig:architecture-flow}
\end{figure}

\subsection{Giao diện chạy thực tế}
\Cref{fig:gui-overview} là ảnh chụp ứng dụng sau khi nhóm em chạy local và nạp bản đồ landmark. Bản đồ hiển thị 65 node và 178 directed road-path edges. Panel bên trái cho phép chọn dataset, thuật toán và scenario. Đây là ảnh chụp từ chương trình đang chạy, không phải hình dựng minh họa.

\begin{figure}[htbp]
\centering
\includegraphics[width=\textwidth]{gui-overview.png}
\caption{Giao diện FloodRoute HCMC khi nạp bản đồ active 65 node.}
\label{fig:gui-overview}
\end{figure}

\subsection{Các luồng chính}
\begin{enumerate}
\item Tìm đường giữa hai điểm: nạp graph, giữ scenario cố định, tính \(\rho_s\), chạy thuật toán và trả về path/cost.
\item So sánh thuật toán: chạy nhiều thuật toán trên cùng graph, start, goal và scenario.
\item Theo dõi trace: lưu frontier, current và closed để giao diện có thể phát lại quá trình tìm kiếm.
\item Tối ưu tour: tạo pairwise matrix rồi chạy Nearest Neighbor hoặc Held-Karp.
\item Mỗi chặng đường đều dùng UCS hoặc A* tùy lựa chọn, trong khi graph ban đầu vẫn được giữ nguyên.
\end{enumerate}

\subsection{Tính tái lập}
Scenario, quy tắc tie-break và phiên bản bộ kết quả được lưu cùng kết quả chạy. Những giá trị mô phỏng được ghi nhãn \texttt{SIMULATED}; những giả định về dữ liệu thực được ghi nhãn \texttt{ASSUMPTION}.
